The statement in the prompt is (slightly) ambiguous about whether repeated summands are allowed.
If repetitions were allowed, the property would be essentially trivial (e.g. a 2-element set
\(\{x,2x\}\) already violates the condition). In the literature around Erd\H{o}s--Straus and Straus,
``admissible'' refers to sums of \emph{distinct} elements.

Accordingly, throughout this section we interpret the problem as follows.

\medskip
\noindent\textbf{Definition (admissible finite set).}
A finite set \(B\subset\mathbb{Z}\) is \emph{admissible} if for any two (finite) \emph{distinct-element} sums
\[
  b_1+\cdots+b_r = b'_1+\cdots+b'_s\qquad (b_i,b'_j\in B,\ b_1<\cdots<b_r,\ b'_1<\cdots<b'_s),
\]
we necessarily have \(r=s\).
Equivalently, if we define for \(r\ge1\)
\[
  S_r(B)\;:=\;\Big\{\, b_1+\cdots+b_r : b_1<\cdots<b_r\in B\,\Big\},
\]
then \(B\) is admissible if and only if \(S_r(B)\cap S_s(B)=\emptyset\) for all \(r\neq s\).

\medskip
\noindent\textbf{Definition of \(h(n)\).}
Let \(h(n)\) be the largest integer such that every set \(A\subset\mathbb{Z}\) with \(|A|=n\)
contains an admissible subset \(B\subseteq A\) with \(|B|\ge h(n)\).
